% 付録再編草案・第3部：数列全体と母関数

\section{数列全体を一つの対象にまとめると,何が見えるか}

ここまでは,次の項を作る規則や,隣り合う項の変化を見てきた.
しかし,一つずつの項ではなく,数列全体の構造をまとめて扱う方法もある.
数列の最初の数項を
\[
  a_0+a_1x+a_2x^2+a_3x^3+\cdots
\]
と並べてみよう.
$x$ は,それぞれの項が何番目にあるかを記録する印である.
この無限和を
\[
  A(x)=\sum_{n=0}^{\infty}a_nx^n
\]
と書いたものが母関数である.
まずは収束する関数としてではなく,数列を係数として並べた記号として考える.

この書き方では,数列の操作が式の操作に移る.
たとえば $A(x)$ に $x$ を掛けると,
\[
  xA(x)=a_0x+a_1x^2+a_2x^3+\cdots
\]
となり,係数を変えずに項を一つずつ右へずらした形になる.
また,
\[
  c_n=\sum_{k=0}^{n}a_kb_{n-k}
\]
という畳み込みは,二つの母関数を掛けたときの $x^n$ の係数になる.
数列のシフトは $x$ 倍に,畳み込みは積に翻訳される.
母関数が便利なのは,無限個の項を圧縮するからだけではなく,数列の操作を代数計算へ移せるからである.

この翻訳をフィボナッチ数列で試す.
\[
  F_0=0,\quad F_1=1,\qquad F_{n+2}=F_{n+1}+F_n
\]
とし,\\,$F(x)=\sum_{n=0}^{\infty}F_nx^n$ と置く.
漸化式を係数ごとに並べると,
\[
  F(x)=x+xF(x)+x^2F(x)
\]
となるから,
\[
  F(x)=\frac{x}{1-x-x^2}
\]
である.
数列を作る規則が,母関数では分母に現れた.
この分母の零点は,特性方程式 $r^2-r-1=0$ の特性根の逆数である.
行列の固有値,特性方程式,母関数の分母は,同じフィボナッチ数列を別々の表現で見たものだった.

母関数は,線形漸化式だけでなく,項どうしの積が現れる漸化式にも使える.
カタラン数は
\[
  C_0=1,\qquad C_{n+1}=\sum_{k=0}^{n}C_kC_{n-k}
\]
で定められる.
右辺は畳み込みだから,\\,$C(x)=\sum_{n=0}^{\infty}C_nx^n$ と置けば,
\[
  C(x)=1+xC(x)^2
\]
となる.
各項を順に計算する漸化式が,母関数についての二次方程式へ変わった.
ここでは母関数が,数列の演算を代数の演算に翻訳したことで,非線形な関係を扱いやすくしている.

ここで,最初に見たフィボナッチ数列へ戻ろう.
特性方程式の根のうち絶対値の大きい方は
\[
  \varphi=\frac{1+\sqrt5}{2}
\]
であり,項の長期的な増加を支配する.
実際,
\[
  F_n\sim\frac{\varphi^n}{\sqrt5}
\]
である.
そのため,正確な一般項を毎回計算しなくても,大きな $n$ での成長を見積もることができる.
母関数の分母からも,同じ特性根が成長を決めることが分かる.
一般項,行列の固有値,母関数,長期的な増大は,別々の話ではなく,一つの更新規則を異なる問いから見た結果である.

漸化式の使われ方は,ここで挙げた話題に限られない.
組合せ論では,最後の選択で場合分けすることで二項係数の漸化式が生まれる.
確率では,状態遷移を行列で表し,長い時間の後の分布を調べる.
数論では,ユークリッドの互除法のように,余りを小さくしながら有限回で止まる計算を記述する.
どの場合も,項を計算するだけでなく,何を状態とし,どの規則で更新し,何を知りたいのかを決めている.

この付録では,本文ですでに扱った不動点や一次分数変換を出発点に,漸化式を状態の更新として読み直した.
更新を線形化すると行列と特性方程式が現れ,その同じ三項間の規則からチェビシェフ多項式が生まれた.
一歩の変化を細かくすると微積分につながり,根へ近づく更新を設計するとニュートン法になった.
数列全体を母関数にまとめると,特性方程式や畳み込みが別の姿を見せた.
話題を一つずつ覚えることが目的ではない.
同じ漸化式を違う言葉で眺めることで,どの構造が見え,どの問題が解けるようになるかを知ることが目的である.
